\paragraph{Inversion.} A soft-binding value is published inside the manifest, so we asked how much of the image it gives away. A decoder trained on the descriptors of the \val{fp:sec/inv/ntrain} distractors reconstructed a $64\times64$ thumbnail of each of \val{fp:sec/inv/ntest} registered images from its descriptor alone (Fig.~\ref{fig:fp-inversion}). Reconstructions from the float descriptors were recognisable in colour, layout and object class (LPIPS \val{fp:sec/inv/DINOv2-S/lpips} for DINOv2-S against \val{fp:sec/inv/mean-image/lpips} for the mean image), and using them as queries retrieved the true image among all \val{fp:sec/inv/ntest} in \val{fp:sec/inv/OpenCLIP-B32/reid/pct}\% (OpenCLIP), \val{fp:sec/inv/DINOv2-S/reid/pct}\% (DINOv2-S) and \val{fp:sec/inv/SSCD-mixup/reid/pct}\% (SSCD) of cases, 40 to 95 times chance. The binary hashes leaked less but still identified images above chance (PDQ \val{fp:sec/inv/PDQ/reid/pct}\%). The keyed projection behaved as intended: an attacker who knows the key still identifies images at about 40 times chance (\val{fp:sec/inv/DINOv2-S-LSH256/reid/pct}\%, half the rate from the float DINOv2-S descriptor), but a decoder trained with any other key reconstructs only the mean image (LPIPS \val{fp:sec/inv/DINOv2-S-LSH256-wrongkey/lpips}) and identifies images at chance (\val{fp:sec/inv/DINOv2-S-LSH256-wrongkey/reid/pct}\%).

\begin{figure*}[!t]
\centering
\includegraphics[width=\textwidth]{fp_fig10_inversion}
\caption{Inverting stored fingerprints. (a) Registered images (top) and thumbnails reconstructed from their descriptors by a decoder trained on distractor descriptors; bottom rows: the keyed 256-bit projection decoded by an attacker with the right and with a wrong key. (b) Share of reconstructions that retrieve their own original among the \val{fp:sec/inv/ntest} registered images (95\% Clopper--Pearson intervals); dashed: chance.}
\label{fig:fp-inversion}
\end{figure*}
